% 電磁気学1 電磁気学1-3 問題1（001-tani。当時は 電磁気学1-1-2 問題1・shou01-02-l1-tani） 解いた日 2026-10-11
%% kotae
\toi{1}{◎}{}
\[ \mathrm{s^{2}\,C^{2}\,kg^{-1}\,m^{-3}}\cdot\mathrm{kg\,m\,C^{-2}}=\mathrm{s^{2}\cdot m^{-2}} \]
\toi{2}{◎}{}
\[ \frac{1}{\sqrt{8.854\times4\times\pi\times10^{-19}}}=\qty{2.998e8}{\metre\per\second}=c \]
\toi{3}{△}{単位}
\begin{align*}
\frac{1}{4\pi\varepsilon_{0}}&=8.99\times10^{9}\\
&\qquad\bigl[\mathrm{s^{-2}\,C^{-2}\,kg\,m^{3}}=\mathrm{N}\cdot\ayamari{\mathrm{m}}{\mathrm{m}^{2}}\,\mathrm{C}^{-2}\bigr]
\end{align*}
（余白に）$8.99\times10^{9}\ \mathrm{N}\,\ayamari{\mathrm{m}}{\mathrm{m}^{2}}\,\mathrm{C}^{-2}$
\naze{$\mathrm{kg\,m^{3}\,s^{-2}\,C^{-2}}=(\mathrm{kg\,m\,s^{-2}})\cdot\mathrm{m^{2}\,C^{-2}}$ なので，正しくは $\qty{8.99e9}{\newton\metre\squared\per\coulomb\squared}$ である。\unit{\newton} に含まれる \unit{\metre} を1つ引いた残りは $\mathrm{m^{2}}$ で，指数が落ちている。数値は正しい。}
%% kizuki
\item 単位を \unit{\newton} を含む形に書き換えるときは，$\unit{\newton}=\mathrm{kg\,m\,s^{-2}}$ を取り出した残りの指数（ここでは $\mathrm{m^{2}}$）を数え直す。クーロンの法則 $F=q^{2}/(4\pi\varepsilon_{0}r^{2})$ から $1/(4\pi\varepsilon_{0})=Fr^{2}/q^{2}$ なので $\mathrm{N\,m^{2}\,C^{-2}}$，と次元で検算できる。
